% ch3_answers.tex -- Module 8, Chapter 3: answers to every Quick check question.
\begin{chapteranswers}{3}{The Quantum Atom}
  % ---- 3.1 Rutherford and Bohr
  \item \textbf{Collapse:} the orbiting electron keeps changing direction, so it is
        accelerating. By Maxwell's theory an accelerating charge gives out EM waves, which carry
        energy away. Its total energy falls, so it moves to smaller and smaller orbits and should
        spiral into the nucleus in about \qty{e-11}{s}. Real atoms are stable.
        \textbf{Spectrum:} as it spirals in, it goes round faster and faster. Classically the
        light has the orbit's frequency, so the frequency would rise smoothly through every
        value: a continuous spectrum. Real hydrogen gives a line spectrum.
  \item $E = E_\mathrm{i} - E_\mathrm{f} = \qty{-0.38}{eV} - (\qty{-1.51}{eV}) = \qty{1.13}{eV}$.
        In joules: $1.13 \times \num{1.602e-19} = \qty{1.81e-19}{J}$.
        Answer: \final{$E = \qty{1.13}{eV} = \qty{1.81e-19}{J}$} (a Paschen line, infrared).
  \item (a) Different atoms can take different routes down, so draw every downward arrow:
        $4\to3$, $4\to2$, $4\to1$, $3\to2$, $3\to1$, $2\to1$.
        Answer: \final{6 different energies} (6 spectral lines).
        (b) The biggest drop gives the highest frequency ($E = hf$): \final{$4 \to 1$}, with
        $E = \qty{-0.85}{eV} - (\qty{-13.6}{eV}) = 12.75 \approx \qty{12.8}{eV}$ (to 1 decimal
        place, like 13.6).
  \item The gaps from $n = 1$ are $\qty{-3.40}{eV} - (\qty{-13.6}{eV}) = \qty{10.2}{eV}$ (to
        $n = 2$) and $\qty{-1.51}{eV} - (\qty{-13.6}{eV}) = \qty{12.1}{eV}$ (to $n = 3$).
        Gaps to higher levels are even bigger (\qty{12.8}{eV} to $n = 4$, and so on, up to
        \qty{13.6}{eV}). So \qty{11.0}{eV}, which lies between \qty{10.2}{eV} and
        \qty{12.1}{eV}, matches no gap. It is also less than the \qty{13.6}{eV} needed to
        ionise the atom. So \final{it is not absorbed}: it passes straight through.
  % ---- 3.2 Rydberg
  \item $\dfrac{1}{\lambda} = R\left(\dfrac{1}{1^2} - \dfrac{1}{2^2}\right)
        = \qty{1.097e7}{m^{-1}} \times 0.75 = \qty{8.228e6}{m^{-1}}$, so
        $\lambda = \dfrac{1}{\qty{8.228e6}{m^{-1}}} = \qty{1.215e-7}{m}$.
        Answer: \final{$\lambda = \qty{122}{nm}$}, ultraviolet (Lyman series).
  \item Put $n_\mathrm{i} = \infty$, so $\dfrac{1}{\infty^2} = 0$:
        $\dfrac{1}{\lambda} = \qty{1.097e7}{m^{-1}} \times \left(\dfrac{1}{4} - 0\right) =
        \qty{2.743e6}{m^{-1}}$, so $\lambda = \qty{3.646e-7}{m}$.
        Answer: \final{$\lambda = \qty{365}{nm}$}. It is \emph{not} visible: it is ultraviolet
        (visible light is about \qtyrange{400}{700}{nm}).
  \item The electron starts at $n = 1$. For absorption the bracket comes out negative, so
        ignore the minus sign and write:
        $\dfrac{1}{\lambda} = R\left(\dfrac{1}{1^2} - \dfrac{1}{n^2}\right)$, so
        $1 - \dfrac{1}{n^2} = \dfrac{1}{R\lambda} =
        \dfrac{1}{\qty{1.097e7}{m^{-1}} \times \qty{1.026e-7}{m}} = 0.8885$.
        Then $\dfrac{1}{n^2} = 1 - 0.8885 = 0.1115$, so $n^2 = 8.97$ and $n = 3.0$.
        Answer: \final{the electrons jump to $n = 3$}.
  % ---- 3.3 de Broglie
  \item $\lambda = \dfrac{h}{mv} = \dfrac{\qty{6.626e-34}{J.s}}{\qty{1.675e-27}{kg} \times
        \qty{2.2e3}{m.s^{-1}}} = \qty{1.80e-10}{m}$.
        Answer: \final{$\lambda = \qty{1.8e-10}{m}$}. This is about the atom spacing, so
        \emph{yes}, a crystal can diffract these neutrons.
  \item $\lambda = \dfrac{h}{mv}$. With the same $v$, a smaller mass gives a longer
        wavelength, so \final{the electron's wavelength is longer}.
        Factor: $\dfrac{m_\mathrm{p}}{m_\mathrm{e}} = \dfrac{\num{1.673e-27}}{\num{9.109e-31}}
        = 1837$, so about \final{\num{1.84e3} times longer}.
  \item $\lambda = \dfrac{h}{mv} = \dfrac{\qty{6.626e-34}{J.s}}{\qty{1200}{kg} \times
        \qty{1.0}{m.s^{-1}}} = \qty{5.5e-37}{m}$.
        It is true that the car has a de Broglie wavelength, and moving slowly does make
        $\lambda$ longer. But \qty{5.5e-37}{m} is far smaller than any gap (even a nucleus is
        about \qty{e-15}{m}). Diffraction is easy to see only when the gap is about as small
        as $\lambda$. A garage door is about \qty{3}{m} wide, so the car cannot diffract.
        Judgement: \final{the claim is wrong}.
  \item Kinetic energy: $K = qV = \qty{1.602e-19}{C} \times \qty{150}{V} = \qty{2.403e-17}{J}$.
        Speed: $v = \sqrt{\dfrac{2K}{m}} = \sqrt{\dfrac{2 \times \qty{2.403e-17}{J}}
        {\qty{9.109e-31}{kg}}} = \qty{7.26e6}{m.s^{-1}}$ (about $0.024c$, so no relativity
        needed). Wavelength: $\lambda = \dfrac{h}{mv} = \dfrac{\qty{6.626e-34}{J.s}}
        {\qty{9.109e-31}{kg} \times \qty{7.264e6}{m.s^{-1}}} = \qty{1.00e-10}{m}$.
        Answers: \final{$v = \qty{7.26e6}{m.s^{-1}}$} and
        \final{$\lambda = \qty{1.00e-10}{m}$}, the size of atom spacings in crystals.
  % ---- 3.4 Bohr limits and Schrodinger
  \item Any three of:\newline
        \textbullet\ the spectra of atoms with more than one electron;\newline
        \textbullet\ why some lines are brighter than others (intensities);\newline
        \textbullet\ fine structure (one line is really several close lines);\newline
        \textbullet\ the Zeeman effect (a line splits in a magnetic field);\newline
        \textbullet\ the Stark effect (a line splits in an electric field).\newline
        (Bohr gives no reason why stationary states do not radiate. That is a limitation, but
        it is not an observation.)
  \item Both split one spectral line into several. The \textbf{Zeeman effect} happens when the
        atoms are in a \emph{magnetic} field. The \textbf{Stark effect} happens when they are in
        an \emph{electric} field. Bohr's model explains neither.
  \item Any two of these differences:\newline
        \textbullet\ The electron is a wave with a wave function, not a particle on a fixed
        orbit.\newline
        \textbullet\ The model gives only the \emph{probability} of where the electron is
        (orbitals), not an exact path.\newline
        \textbullet\ Quantised energy comes out of the wave equation; it is not assumed.\newline
        \textbullet\ It works (approximately) for atoms with many electrons, not only
        hydrogen.
  \item (a) For $n = 2$: $\lambda = \dfrac{2L}{2}$, half of the $n = 1$ value. So $v$ is 2
        times bigger and $K = \frac12 mv^2$ is $2^2 = 4$ times bigger:
        $K_2 = 4 \times \qty{2.41e-19}{J} = \qty{9.64e-19}{J}$.
        Answer: \final{$K_2 = \qty{9.64e-19}{J}$}.\newline
        (b) Energy is conserved, so the photon carries the energy lost:
        $E = K_2 - K_1 = \qty{9.64e-19}{J} - \qty{2.41e-19}{J} = \qty{7.23e-19}{J}$.
        In eV: $\dfrac{\qty{7.23e-19}{J}}{\num{1.602e-19}} = \qty{4.51}{eV}$.
        Wavelength: $\lambda = \dfrac{hc}{E} = \dfrac{\qty{6.626e-34}{J.s} \times
        \qty{3.00e8}{m.s^{-1}}}{\qty{7.23e-19}{J}} = \qty{2.75e-7}{m} = \qty{275}{nm}$.
        Answer: \final{$E = \qty{4.51}{eV}$, $\lambda = \qty{275}{nm}$}, ultraviolet
        (shorter than \qty{400}{nm}).
        Check: \qty{4.51}{eV} is more than violet light's \qty{3.1}{eV}, so the photon must
        be ultraviolet \cmark.
\end{chapteranswers}
